\section{The Concluding Rites}

\subsection{Four elements, briefly}

The Instruction lists what belongs to the Concluding Rites: brief announcements, if they are
necessary; the greeting and blessing of the priest, which on certain days and occasions is
enriched and expressed by a prayer over the people or another more solemn formula; the dismissal
of the people by the deacon or the priest, so that each may return to his own good works,
praising and blessing God; and the kissing of the altar by priest and deacon, and then the
profound bow toward the altar by priest, deacon, and the other
ministers.\endnote{\latin{IGMR} 90 a--d. The purpose clause in (c) --- ``ut unusquisque ad opera
sua bona revertatur, collaudans et benedicens Deum'' --- is the Instruction's own statement of
what the dismissal is for and is the basis of the reading offered below.} The Order of Mass at
nn.\ 140--146 supplies the texts and the rubrics, including the pontifical form of the blessing
with its versicles \latin{Sit nomen Domini benedictum} and \latin{Adiutorium nostrum in nomine
Domini}, and the provision that if another liturgical action follows immediately the rites of
dismissal are omitted.

Of these, the announcements are conditional and the Instruction's ``if they are necessary'' is a
restriction that the ordinary practice of Sunday parishes tests severely. The blessing is
invariable in its simple form; the solemn blessings and prayers over the people that follow the
Order of Mass in the book are a substantial restored corpus, and the 2008 reprint corrected the
rubric that introduces them, replacing \latin{dicit invitatorium: Inclinate vos ad benedictionem
vel aliis verbis expressum} with \latin{dicit invitationem: Inclinate vos ad
benedictionem.}\endnote{\latin{Variationes} in \emph{Notitiae} 44 (2008) 372, the entry for p.\
616. The change does two things: it corrects the term (\latin{invitatio} rather than
\latin{invitatorium}, which properly names the opening psalm of the Office) and it removes the
permission to express the invitation in other words. The 2002 reading was confirmed at the
corresponding place of the registered reproduction, which prints the longer form with
\latin{vel aliis verbis expressum}.} That is a small tightening with a visible effect: after
2008 the invitation to bow for a solemn blessing is a fixed text.

\subsection{The dismissal and its 2008 enlargement}

The 2002 Missal printed one dismissal: \latin{Ite, missa est}, answered \latin{Deo gratias}. The
2008 reprint added three alternatives at n.\ 144, printed in the \latin{Supplementum} under the
heading \latin{Ad n.\ 144 Ordinis Missae}:\endnote{\latin{Variationes} in \emph{Notitiae} 44 (2008) 372 and the \latin{Supplementum} at
printed p.\ 384 of the same issue, both read publication-locally 2026-07-25. The same
\latin{Supplementum} adds, for the Mass of the Vigil of Pentecost, \latin{Vel: Ite in pace,
alleluia, alleluia} beside the existing \latin{Ite, missa est, alleluia, alleluia}. The addition
of alternative dismissal formulas is one of the three matters submitted for papal approval under
Decree N.\ 652/08/L of 8 June 2008, as the 2008 notice records.}


\begin{studyframe}
\latin{Ite, ad Evangelium Domini annuntiandum.}\\
\latin{Vel: Ite in pace, glorificando vita vestra Dominum.}\\
\latin{Vel: Ite in pace.}
\end{studyframe}

This is the most substantive addition the 2008 reprint made to the Order of Mass, and it is a
deliberate act of interpretation. \latin{Ite, missa est} is famously opaque: the noun
\latin{missa} in late Latin means a dismissal, and the formula's literal sense is closer to ``go,
it is the dismissal'' than to anything theological; the word later gave the whole rite its name.
The three added formulas do not translate it. They say what the dismissal is for, in the terms
the Instruction had already used at n.\ 90 c: going to announce the Gospel of the Lord, going in
peace glorifying the Lord by one's life. In other words, the 2008 book supplies, as options
beside the ancient formula, explicit missionary readings of an act whose ancient formula does not
carry them.

Whether that is a gain is a fair question and this study will not pretend it has no cost. The
gain is that the connection the Instruction asserts --- dismissal as sending to good works and
the praise of God --- is now audible in the rite itself, and that the exposure of \latin{missa
est} to a missionary rereading is not an invention of preachers but a provision of the book. The
cost is that an opaque formula whose opacity was itself instructive, and which named the whole
rite, now competes with three transparent ones; and transparency is not always the liturgical
virtue. The book keeps \latin{Ite, missa est} first among the four, which suggests the same
judgement.

The 2011 English of the four dismissals is not reproduced here. The English of \latin{Ite, missa
est} in particular is a translator's problem with no clean solution, since the formula is
untranslatable in the strict sense; the approved English is under copyright and this study does
not print it or paraphrase it closely.

\subsection{The last gesture}

The rite ends as it began: the priest venerates the altar with a kiss as at the beginning, and,
after a profound bow with the ministers, withdraws.\endnote{\latin{Ordo Missae} 145--146.} The
symmetry is exact and it is the Order's own last statement about what has happened. The altar
that was kissed before the assembly was greeted is kissed again before the assembly is left.
Between those two kisses the Church has heard the word, offered the sacrifice, and been fed;
the frame closes on the stone, not on the people, and the last movement of the celebrant's body
is the same act of homage as the first.
